% !TEX root = ../../../tree_notes.tex
% =====================================================================
%  14.3  Greedy on trees: exchange arguments   sec:c27-greedy
%  Exports: def:c27-leaf-prune  thm:c27-exchange  lst:c27-greedy
%           fig:c27-greedyfig  ex:c27-greedy-run  prob:c27-greedy-solved
%           lst:c27-1029e-soln
%  When a budgeted tree DP collapses to a single leaf-pruning sweep because
%  a tree has no cycles: a cut-then-swap exchange does the work of the whole
%  table.  Anchor CF 1029E (Tree with Small Distances): add fewest edges so
%  every vertex is within distance 2 of vertex 1.  Greedy: repeatedly take a
%  DEEPEST not-yet-covered vertex v and attach vertex 1 to parent(v).
%  Verified: samples -> 2, 0, 1; on T the greedy adds (1,9),(1,5) = 2.
%  NOTE: CF 1914F was moved to the Mixed Set -- its "greedy" is a
%  path-invariant/DP hybrid, not a clean exchange argument (see 99_chapter_end).
% =====================================================================
\section{Greedy on Trees: the Exchange Argument}{\label{sec:c27-greedy}}

\begin{previouslybox}
\begin{enumerate}[nosep, leftmargin=1.6em]
  \item \secref{c27-knapsack} merged children under a \emph{budget}. What would
        the state look like if the budget were just ``covered within distance
        $2$ of the root'' -- a Boolean, essentially? (a DP still \emph{works};
        the question is whether a greedy matches it for free)
  \item Ch.~12 found the diameter by ``the farthest node from any node is an
        endpoint'' -- a \emph{greedy} proved by exchange (\thmref{c09-dpjoin}'s
        sibling). What made the swap argument legal on a tree that fails on a
        general graph? (the unique simple path -- swapping a segment cannot
        create a shortcut you did not already account for)
  \item In the rooted tree of Ch.~4, ``deepest uncovered vertex'' is well
        defined. Why is a \emph{deepest} vertex necessarily a leaf of the
        not-yet-covered set? (every child is deeper, hence already processed or
        covered)
  \item Name the greedy invariant you must maintain so the sweep never
        re-counts an edge. (a \lstinline{covered[]} flag updated by a bounded
        local BFS whenever an edge is added)
\end{enumerate}
\end{previouslybox}

\begin{topicintro}
Not every budgeted tree problem wants the array. When the constraint is a
\emph{distance from a fixed root} -- ``bring every vertex within $2$ hops of
vertex~$1$ by adding as few edges as you like'' -- a single top-down tree DP
with a per-node budget would work, but it is overkill: the tree's acyclicity lets
you prove a bare \emph{leaf-pruning greedy} optimal by a cut-then-swap exchange.
This section shows the paradigm: identify a deepest vertex that still violates the
condition, pay for it at the \emph{latest} possible moment (attach near its
parent), and argue that any optimal solution can be edited to make the same choice
without cost. CF~1029E (\probref{c27-greedy-solved}) is the case study, and the
method generalises to every ``cover/guard the tree with the fewest additions''
problem that has a \emph{radius} measured from a fixed centre.

\textbf{Prerequisites.} Ch.~4's rooting by a parent parameter and Ch.~12's
depth (\coderef{c09-dfs3}); \secref{c27-knapsack} for the DP you are \emph{not}
writing and why a greedy can replace it; Ch.~12's exchange intuition behind
\thmref{c09-dpjoin} (a sibling argument, cited not re-proved).

\textbf{After this section you can}
\begin{itemize}[nosep]
  \item define the deepest-uncovered leaf-pruning order and its coverage rule; \hfill\lv{2}
  \item prove a tree greedy optimal with a cut-then-swap exchange; \hfill\lv{4}
  \item cold-write CF~1029E's ``attach $1$ to the parent of the deepest uncovered
        leaf'' sweep and trace it. \hfill\lv{3}--\lv{4}
\end{itemize}
\end{topicintro}

\subsection{Intuition}

Distance-from-a-root constraints are \emph{local upward}: a vertex $v$ at depth
$d > 2$ can only be brought within distance~$2$ of the root by an added edge that
lands \emph{near} $v$ (within one hop of it), because a new edge $(1,u)$ makes the
distance to $v$ equal $1 + \dist(u,v)$, and we need $1 + \dist(u,v) \le 2$, i.e.\
$\dist(u,v) \le 1$. So $u$ must be $v$ itself or a neighbour of $v$. The
\term{exchange} insight is that among $\{v\} \cup \children(v)\cup\{\parent(v)\}$,
the choice $u = \parent(v)$ is never worse: attaching $1$ to $\parent(v)$ covers
$v$ \emph{and} every sibling of $v$ \emph{and} the grandparent (all within one hop
of $\parent(v)$), whereas attaching to $v$ covers only $v$ and $\parent(v)$. Hence
there is an optimal solution that, for the deepest uncovered vertex, uses the edge
to its parent -- the greedy choice -- so greedy is optimal.

The tree's acyclicity is doing quiet work. On a general graph a ``deepest'' vertex
is meaningless (no root, no unique path) and swapping an edge can create shortcuts
that invalidate the coverage accounting; here every vertex has one parent, the path
to the root is unique, and processing \emph{deepest first} guarantees a vertex's
subtree is already settled when you decide about it -- the same post-order discipline
as \secref{c27-knapsack}, but you pay $O(1)$ per vertex instead of running a table.

\subsection{Definitions and Notation}

\begin{definition}{Deepest-uncovered leaf-pruning order}{c27-leaf-prune}
Root the tree at the centre $c$ (here $c = 1$) and compute depths. A vertex is
\term{covered} if its current distance to $c$ is at most~$2$ (with all edges added
so far). The \term{leaf-pruning order} is vertices sorted by \emph{non-increasing
depth}; sweeping it and, at each still-uncovered vertex $v$, taking the action
``add edge $(c, \parent(v))$ and mark $\parent(v)$ and its neighbours covered''
is the greedy. Because a deepest uncovered vertex has no uncovered child (children
are deeper and precede it), every decision is made once, bottom-up.
\end{definition}

\begin{theorem}{The deepest-uncovered greedy is optimal for the distance-$R$ cover}{c27-exchange}
For $R = 2$ and centre~$1$, the greedy of \defref{c27-leaf-prune} adds the minimum
number of edges to make every vertex within distance~$2$ of~$1$.
\end{theorem}
\begin{proofbox}[\thmref{c27-exchange}: exchange a deepest vertex's covering edge onto its parent]
Let $v$ be the deepest vertex uncovered by the initial tree (depth $d \ge 3$), so
$v$ is a leaf among the uncovered. In \emph{any} valid solution $O$, $v$ ends within
distance~$2$ of~$1$. Since $\dist(1,v) = d \ge 3$ originally, $O$ must contain an
added edge $(1,u)$ with $1 + \dist(u,v) \le 2$, i.e.\ $\dist(u,v) \le 1$: $u \in
\{v\} \cup \mathrm{nbr}(v)$. As $v$ is a leaf (no children), $\mathrm{nbr}(v) = \{\parent(v)\}$,
so $u \in \{v, \parent(v)\}$. We claim swapping that edge to $(1,\parent(v))$ keeps
$O$ valid and no larger. The coverage of an added edge $(1,x)$ is the closed
neighbourhood $\{x\} \cup \mathrm{nbr}(x)$ measured in new distance $1 + \dist(x,\cdot)
\le 2$, i.e.\ vertices within distance~$1$ of $x$. For $x = \parent(v)$ this set is
$\{\parent(v), \text{grandparent}, \text{all of } v\text{'s siblings}, v\}$,
which contains $\{v, \parent(v)\}$ (the set for $x = v$). So replacing $(1,v)$ by
$(1,\parent(v))$ covers a superset: every vertex still within distance~$2$ of~$1$,
and $|O|$ is unchanged. Hence there is an optimal solution $O'$ \emph{containing}
the greedy edge $(1,\parent(v))$. Both the greedy and $O'$ now cover $v$ and
everything $(1,v)$ would have, so remove $v$'s branch and induct on the remaining,
strictly shallower tree: the greedy is optimal. \qed
\end{proofbox}

\begin{insightbox}
The greedy's power is exactly the \emph{superset} relation
$\mathrm{cov}(\parent(v)) \supseteq \mathrm{cov}(v)$. That is a property of rooted
trees (a parent's closed neighbourhood swallows the child's when the child is a
leaf), not of general graphs. It is the same reason Ch.~12's ``farthest node is a
diameter endpoint'' swap worked and the analogous claim fails on a cyclic graph:
the unique path keeps the coverage accounting honest. When you see a distance-from-
a-fixed-centre objective, \emph{test} for this superset before building a DP -- if
it holds, the greedy is free.
\end{insightbox}

\subsection{The Algorithm: Deepest-Uncovered Sweep}

\begin{lstlisting}[caption={The CF 1029E greedy: root at 1, sweep vertices by decreasing depth, attach 1 to the parent of each still-uncovered vertex}, label={lst:c27-greedy}]
// Goal: every vertex within distance 2 of vertex 1; add fewest edges.
// root = 1; par[v], dep[v] from one DFS; covered[v] = "dist to 1 <= 2 so far".
int R = 2;                             // radius target
vector<vector<int>> adj;
vector<int> par, dep, ord;
vector<char> covered;

void dfs(int v, int p) {
    par[v] = p;
    dep[v] = (p == 0 ? 0 : dep[p] + 1);
    for (int u : adj[v]) if (u != p) dfs(u, v);
}

int greedy() {
    dfs(1, 0);
    for (int i = 1; i <= (int)adj.size() - 1; i++) covered[i] = (dep[i] <= R);
    iota(ord.begin(), ord.end(), 1);   // 1..n
    sort(ord.begin(), ord.end(), [&](int a, int b){ return dep[a] > dep[b]; });
    int ans = 0;
    for (int v : ord) {
        if (covered[v]) continue;      // already within distance R of 1
        int p = par[v];                // attach at the PARENT (exchange proof)
        ans++;
        covered[p] = 1;                // p and its neighbours now reach 1 in <= 2
        for (int u : adj[p]) covered[u] = 1;
    }
    return ans;
}
\end{lstlisting}

\begin{complexitybox}
One DFS for depth/parent $\bigO{n}$; the sort is $\bigO{n\log n}$ (a counting sort on
depth makes it $\bigO{n}$ since depths are in $[0,n]$); the sweep marks each vertex
covered once but the neighbour loop re-reads edges, so the total marking is
$\bigO{\sum_v \deg(v)} = \bigO{n}$. Overall $\bigO{n \log n}$ dominated by the sort.
Memory $\bigO{n}$. Verdict: a budgeted DP would be $\bigO{nR}$ with a state array;
the greedy is the same order but with a \emph{proof} (exchange) that a single number
per vertex suffices. Watch recursion depth $\bigO{h}$ on chains ($n \le 2\cdot10^5$).
\end{complexitybox}

\subsection{Worked Example: The Two-Edge Cover of $T$}

\begin{example}{Running CF~1029E's greedy on the running tree $T$}{c27-greedy-run}
Root $T$ at $1$; depths are $(0,1,1,2,2,2,2,3,3,4)$ for $v = 1..10$, so the vertices
deeper than~$2$ are $8$ (depth~$3$), $9$ (depth~$3$) and $10$ (depth~$4$). Sweep by
decreasing depth and attach $1$ to the parent of each uncovered vertex. The
result is \emph{two} added edges, and \thmref{c27-exchange} says two is optimal.
\end{example}

\begin{figure}[H]
\centering
\begin{tikzpicture}[diagram, scale=0.82, transform shape,
    tnode/.style={circle, draw=primaryblue, line width=0.9pt, fill=primaryblue!8,
                  minimum size=7mm, font=\small\bfseries, text=primaryblue!60!black}]
  \node[tnode] (n1) at (0,0) {$1$};
  \node[tnode] (n2) at (-3,-1.6) {$2$};
  \node[tnode] (n3) at ( 3,-1.6) {$3$};
  \node[tnode] (n4) at (-5.2,-3.2) {$4$};
  \node[tnode] (n5) at (-1.6,-3.2) {$5$};
  \node[tnode] (n6) at ( 1.6,-3.2) {$6$};
  \node[tnode] (n7) at ( 5.2,-3.2) {$7$};
  \node[tnode] (n8) at (-1.6,-4.8) {$8$};
  \node[tnode] (n9) at ( 5.2,-4.8) {$9$};
  \node[tnode] (n10) at ( 5.2,-6.4) {$10$};
  \draw[tedge] (n1)--(n2) (n1)--(n3) (n2)--(n4) (n2)--(n5)
               (n3)--(n6) (n3)--(n7) (n5)--(n8) (n7)--(n9) (n9)--(n10);
  % added edges 1-9 and 1-5
  \draw[tedgehl, line width=1.8pt, dashed] (n1)--(n9);
  \draw[tedgehl, line width=1.8pt, dashed] (n1)--(n5);
  \node[font=\scriptsize, text=warningred!75!black, anchor=east] at (-0.2,-2.6) {add (1,5)};
  \node[font=\scriptsize, text=warningred!75!black, anchor=west] at (4.4,-1.0) {add (1,9)};
  % depth badges
  \foreach \nd/\dd in {n1/0,n2/1,n3/1,n4/2,n5/2,n6/2,n7/2,n8/3,n9/3,n10/4} {
    \node[font=\tiny\ttfamily, text=softpurple, anchor=north west] at (\nd.north east) {\dd};}
  \node[annot, font=\scriptsize, text=difforange!75!black, anchor=west] at (7.0,-3.2)
    {dashed red = added};
  \node[annot, font=\scriptsize, text=softpurple, anchor=west] at (7.0,-3.9)
    {gray = depth (badges)};
  \node[annot, font=\scriptsize, text=successgreen!50!black, anchor=west] at (7.0,-4.8)
    {answer $=2$ edges};
\end{tikzpicture}
\caption{The greedy cover of $T$. Deepest uncovered node $10$ (depth~$4$) forces the
edge $(1,\parent(10)) = (1,9)$, covering $\{9,7,10\}$ within distance~$2$ of~$1$; then
the deepest still-uncovered node $8$ (depth~$3$) forces $(1,\parent(8)) = (1,5)$,
covering $\{5,2,8\}$. Every vertex is now within distance~$2$ of~$1$, and
\thmref{c27-exchange} guarantees no single-edge alternative does it (one edge can
cover at most a closed neighbourhood, which never spans both the $8$- and $10$-branches).
Answer $=2$.}
\label{fig:c27-greedyfig}
\end{figure}

\begin{tracebox}[The leaf-pruning sweep on $T$, decreasing depth]
covered$=$depth~$\le2$: $\{1,2,3,4,5,6,7\}$. Uncovered: $10(4), 8(3), 9(3)$.\\
order by depth desc: $10, 8, 9, \dots$\\
$10$: uncovered. $p=\parent(10)=9$. add $(1,9)$, ans$=1$. mark $9$ and
$\mathrm{nbr}(9)=\{7,10\}$ covered.\\
$8$: uncovered. $p=\parent(8)=5$. add $(1,5)$, ans$=2$. mark $5$ and
$\mathrm{nbr}(5)=\{2,8\}$ covered.\\
$9$: covered (skip). all of $8,9,10$ now covered. \;\;Sweep ends, ans$=\mathbf{2}$. \checkmark
\end{tracebox}

\begin{insightbox}[Key takeaways]
\begin{enumerate}[nosep, leftmargin=1.6em]
  \item \emph{Attach at the parent, not the vertex.} $(1,\parent(v))$ covers a
        superset of $(1,v)$; that single fact is the entire exchange.
  \item \emph{Deepest-first is a post-order.} Processing by non-increasing depth means
        a vertex's whole (deeper) subtree is decided before the vertex -- no
        back-tracking, one decision each.
  \item \emph{Coverage marking is the invariant.} The answer is only correct because
        you re-mark $\parent(v)$ and its neighbours; skipping the re-mark is the classic
        off-by-one (Trap~1 below).
\end{enumerate}
\end{insightbox}

\subsection{Variations and Twists}

\begin{itemize}[nosep]
  \item \textbf{Target radius $R$ general}: for distance-$R$ cover the exchange reads
        ``attach at $v$'s $(R{-}1)$-st ancestor'' and mark that ancestor's radius-$(R{-}1)$
        ball; the greedy is still optimal by the same superset argument. \hfill\lv{4}
  \item \textbf{Guarding, not covering}: ``minimum vertex set so every edge touches
        one'' (vertex cover on a tree) is \emph{not} a distance-from-centre problem --
        its greedy needs the in/out pair of Ch.~13 (\secref{c10-twostate}), a warning
        about over-applying the leaf-prune heuristic. \hfill\lv{4}
  \item \textbf{When the DP wins}: if the per-vertex cost varies (each added edge has a
        price), the superset relation breaks and you \emph{do} need a budgeted state --
        hand back to \secref{c27-knapsack}. \hfill\lv{4}
  \item \textbf{Centroid not fixed}: CF~1029E fixes the centre at~$1$; asking for the best
        centre is the radius problem, Ch.~12's \secref{c09-radius}, and rerooting it is
        Ch.~15 (\chref{c20}). \hfill\lv{5}
\end{itemize}

\subsection{Common Pitfalls}

\begin{warningbox}[Trap 1 -- attach to the vertex, not its parent]
Connecting $1$ directly to the deepest leaf $v$ fails the exchange: $\mathrm{cov}(v)$
misses $v$'s siblings, so you spend extra edges covering them one by one and return a
non-minimal count. Symptom: correct on chains (where a leaf has no siblings) but too
high on stars and brooms. Attach at $\parent(v)$.
\end{warningbox}

\begin{warningbox}[Trap 2 -- marking only $v$, not $\parent(v)$'s whole neighbourhood]
After adding $(1,\parent(v))$ you must mark $\parent(v)$ \emph{and every neighbour of}
$\parent(v)$ as covered -- the grandparent and all siblings too. Marking only $v$ lets
a sibling re-trigger a redundant edge. The coverage rule is ``vertices at distance
$\le 1$ from $\parent(v)$'', i.e.\ the closed neighbourhood.
\end{warningbox}

\begin{warningbox}[Trap 3 -- a wrong ``deepest'' (using index, not depth)]
The sweep must be ordered by \emph{depth}, not vertex number or input order; and among
equal depths the order is arbitrary but must be processed \emph{after} all deeper
vertices. Sorting by index silently breaks the post-order guarantee and can attach an
edge at a vertex whose parent is itself not yet finalised.
\end{warningbox}

\subsection{Problem Lab}

\solvedproblem{c27-greedy-solved}{Tree with Small Distances}{Codeforces 1029E}
  {https://codeforces.com/problemset/problem/1029/E}{CF 1800 · L4}
  {leaf-pruning greedy, exchange argument, distance-from-centre cover,
   closed-neighbourhood marking, depth sort}

\begin{problemstatement}[Tree with Small Distances (statement and samples verified against the live problem)]
Given a tree of $n$ vertices ($2 \le n \le 2\cdot10^5$). Add the \emph{minimum} number
of edges so that the shortest-path distance from vertex~$1$ to every other vertex is
at most~$2$. No loops or multiple edges may be added. Output that minimum number.
\textbf{Example~1}: $n=7$, edges $1\text{--}2,2\text{--}3,2\text{--}4,4\text{--}5,
4\text{--}6,5\text{--}7$; output \texttt{2} (e.g.\ add $(1,5)$ and $(1,6)$).
\textbf{Example~2}: $n=7$, edges $1\text{--}2,1\text{--}3,2\text{--}4,2\text{--}5,
3\text{--}6,1\text{--}7$; output \texttt{0} (every vertex is already within distance~$2$
of~$1$). \textbf{Example~3}: $n=7$, edges $1\text{--}2,2\text{--}3,3\text{--}4,
3\text{--}5,3\text{--}6,3\text{--}7$; output \texttt{1} (the only optimal addition is
$(1,3)$).
\end{problemstatement}

\begin{attemptgate}[Attempt it yourself first -- \lv{3}, give yourself 20 minutes]
Resist the DP. Ask: to bring a deep vertex $v$ within distance~$2$ of~$1$ with one new
edge, where can that edge's far endpoint be? Then which choice covers the most.
\peek{Adding edge $(1,u)$ sets the distance to $v$ as $1 + \dist(u,v)$. For that to be
      $\le 2$, how far can $u$ be from $v$?}
\nudge{$\dist(u,v) \le 1$: so $u$ is $v$ or a neighbour. For a deepest leaf $v$, its
      only neighbour is $\parent(v)$ -- and the parent also covers $v$'s siblings.}
\fullhint{Root at~$1$, compute depths. Mark every vertex of depth $\le 2$ covered.
Process remaining vertices in decreasing depth: if $v$ is still uncovered, add edge
$(1,\parent(v))$, increment the answer, and mark $\parent(v)$ together with all of
$\parent(v)$'s neighbours covered. This is optimal by \thmref{c27-exchange}
(the parent-edge dominates the vertex-edge). Answer $=$ number of added edges.}
\end{attemptgate}

\solstep{Understanding the Problem}
Only distance from the fixed centre $1$ matters, and an added edge $(1,u)$ acts locally:
it puts $u$ at distance~$1$ and $u$'s neighbours at distance~$2$, and nothing farther.
So a single edge ``fixes'' exactly a closed neighbourhood $\{u\}\cup\mathrm{nbr}(u)$. Vertices
of depth $\le 2$ need no help. For a deeper vertex $v$, the cheapest way to include it in
some closed neighbourhood is to centre that neighbourhood on $\parent(v)$, which also
absorbs $v$'s siblings and grandparent. The objective (minimise added edges) then becomes
a covering problem on the tree's lower part, and the nesting of closed neighbourhoods in a
rooted tree is exactly what the exchange turns into a greedy.

\solstep{Key Observations}
\begin{observation}
An added edge $(1,u)$ covers precisely the vertices at distance $\le 1$ from $u$
(their new distance to~$1$ is $1+\dist(u,\cdot)\le 2$). Coverage is a \emph{closed
neighbourhood}, nothing more.
\end{observation}
\begin{observation}
For a deepest uncovered vertex $v$ (depth $\ge 3$), the only edges that can cover $v$
are $(1,v)$ or $(1,\parent(v))$; and $\mathrm{cov}(\parent(v)) \supseteq \mathrm{cov}(v)
\cup\{\text{siblings},\text{grandparent}\}$. Choosing the parent is never worse -- the
exchange of \thmref{c27-exchange}.
\end{observation}
\begin{observation}
Ordering by decreasing depth makes every decision final: a vertex's descendants are
strictly deeper and already covered, so the vertex is a leaf of the uncovered set when
processed -- the post-order property that keeps the greedy single-pass.
\end{observation}

\solstep{Think It Through With Me}
\begin{thinkbox}
\deadend{Run a DP $\mathit{dp}[v][0/1/2]$ = ``$v$ is within $0$, $1$, or $2$ hops of a
connected-to-$1$ node'' and merge children like \secref{c27-knapsack}.}
{Correct, but you are paying an array-and-merge price for a problem whose state is
essentially ``is covered yet''. Once you notice the closed-neighbourhood superset, the
whole table collapses to one flag. The DP here is a solution looking for a problem.}
\deadend{Greedy: connect $1$ to the deepest leaf $v$ itself.}
{Covering $v$ but not its siblings -- on Example~1 you would spend an edge per leaf of
a broom and over-count. The fix is to attach one level up, at $\parent(v)$: same proof,
strictly more coverage. Trap~1 made concrete.}
\aha{Sweep by \emph{decreasing depth}; at each uncovered $v$ add $(1,\parent(v))$ and
mark $\parent(v)$'s closed neighbourhood covered. Optimal by the \thmref{c27-exchange}
superset, single pass because deepest-first is post-order.}
\end{thinkbox}

\solstep{Solution}
DFS from~$1$ for $\parent(\cdot)$ and $\depth(\cdot)$; mark $\mathit{covered}[v]$ true when
$\depth(v) \le 2$. Sort vertices by decreasing depth. For each uncovered $v$: answer++,
let $p = \parent(v)$, set $\mathit{covered}[p] = $ true and $\mathit{covered}[u] = $ true for every
$u \in \mathrm{nbr}(p)$. Print the answer.

\solstep{Proof of Correctness}
\begin{proofbox}[\coderef{c27-greedy} adds a minimum edge set making all distances to~$1$ at most $2$]
\emph{Validity.} Each added edge is $(1,\parent(v))$ for a distinct uncovered $v$ (once
covered, $v$ is skipped, so no edge is added twice for it), and none is a loop or
duplicate ($\parent(v) \ne 1$ because $v$ has depth $\ge 3$, and $(1,\parent(v))$ was not
an original edge since $\parent(v)$ has depth $\ge 2$ while all depth-$\ge2$ vertices
adjacent to~$1$ would have depth~$1$). After the sweep every vertex is $\mathit{covered}$, i.e.\
within distance~$2$ of~$1$ in the augmented graph -- by construction $\mathit{covered}$ is set
exactly when some added edge's closed neighbourhood includes it or its depth~$\le2$.
\emph{Minimality.} By \thmref{c27-exchange}, for the deepest uncovered $v$ some optimal
solution contains the greedy edge $(1,\parent(v))$; contract the covered neighbourhood and
apply the same argument inductively on strictly shallower residual trees. Each greedy step
matches one unavoidable edge, so the greedy count equals the optimum. \qed
\end{proofbox}

\solstep{Implementation}
\begin{lstlisting}[caption={Codeforces 1029E -- deepest-uncovered leaf-pruning greedy (\coderef{c27-greedy}), parent-attachment}, label={lst:c27-1029e-soln}]
#include <bits/stdc++.h>
using namespace std;
int n;
vector<vector<int>> adj;
vector<int> par, dep, ord;
vector<char> cov;

void dfs(int v, int p) {
    par[v] = p;
    dep[v] = (p == 0 ? 0 : dep[p] + 1);
    for (int u : adj[v]) if (u != p) dfs(u, v);
}

int main() {
    if (scanf("%d", &n) != 1) return 0;
    adj.assign(n + 1, {});
    for (int i = 0; i < n - 1; i++) {
        int u, v; scanf("%d %d", &u, &v);
        adj[u].push_back(v); adj[v].push_back(u);
    }
    par.assign(n + 1, 0); dep.assign(n + 1, 0); cov.assign(n + 1, 0);
    ord.assign(n, 0);
    dfs(1, 0);
    for (int i = 1; i <= n; i++) cov[i] = (dep[i] <= 2);   // already within 2 of 1
    iota(ord.begin(), ord.end(), 1);
    sort(ord.begin(), ord.end(), [&](int a, int b){ return dep[a] > dep[b]; });
    int ans = 0;
    for (int v : ord) {
        if (cov[v]) continue;
        int p = par[v];              // attach at the parent -- exchange proof
        ans++;
        cov[p] = 1;                  // p's closed neighbourhood now reaches 1 in <= 2
        for (int u : adj[p]) cov[u] = 1;
    }
    printf("%d\n", ans);
    return 0;
}
\end{lstlisting}

\solstep{Trace}
\begin{tracebox}[Official Example~1: $n=7$, edges $1\text{--}2,2\text{--}3,2\text{--}4,4\text{--}5,4\text{--}6,5\text{--}7$ $\Rightarrow$ 2]
Parents: $\parent(2)=1,\parent(3)=2,\parent(4)=2,\parent(5)=4,\parent(6)=4,\parent(7)=5$.
Depths: $1\!:\!0,2\!:\!1,3\!:\!2,4\!:\!2,5\!:\!3,6\!:\!3,7\!:\!4$. Initially covered
(depth~$\le2$): $\{1,2,3,4\}$. Decreasing-depth order: $7(4),5(3),6(3),3,4,2,1$.\\
$7$: uncovered. $p=\parent(7)=5$. add $(1,5)$, ans$=1$. cover $5$ and
$\mathrm{nbr}(5)=\{4,7\}$ $\Rightarrow$ $\{5,4,7\}$ covered.\\
$5$: covered (skip). \;\;$6$: uncovered. $p=\parent(6)=4$. add $(1,4)$? no --
$\parent(6)=4$, and $4$ is an original neighbour of $2$, but is $(1,4)$ a legal new
edge? $4$ has depth~$2$, not adjacent to~$1$ in the input, so $(1,4)$ is legal. ans$=2$.
cover $4$ and $\mathrm{nbr}(4)=\{2,5,6\}$ $\Rightarrow$ $6$ covered.\\
Rest covered. Total added $=\mathbf{2}$. \checkmark\ (CF lists $(1,5),(1,6)$ as an
equivalent optimum; $(1,5),(1,4)$ covers the same set.)
\end{tracebox}

\begin{tracebox}[Examples 2 and 3]
Ex.~2: all depths $\le 2$ ($4,5,6$ at depth~$2$, $2,3,7$ at depth~$1$) $\Rightarrow$ no
uncovered vertex, ans$=0$. \checkmark\\
Ex.~3: $\parent(2)=1,\parent(3)=2,\parent(4)=\parent(5)=\parent(6)=\parent(7)=3$;
depths $2,3$ for $3$, and $4,5,6,7$ at depth~$3$. Only depth-$\ge3$ are $4,5,6,7$.
First (say $4$): $p=\parent(4)=3$, add $(1,3)$, cover $3$ and $\mathrm{nbr}(3)=\{2,4,5,6,7\}$
$\Rightarrow$ all four leaves covered at once. Remaining $5,6,7$ skipped. ans$=\mathbf{1}$.
\checkmark\ (CF: the only optimal edge is $(1,3)$.)
\end{tracebox}

\solstep{Complexity}
\begin{complexitybox}
$\bigO{n\log n}$ (depth sort) or $\bigO{n}$ with counting sort, one DFS, one sweep, each
edge scanned at most a constant number of times in neighbourhood marking. Memory $\bigO{n}$.
$n\le2\cdot10^5$ so the $\bigO{n\log n}$ sort is trivially fine; guard the DFS recursion on
a chain (\re risk) with the two-phase stack. Verdict: \ac; the lesson is not the code but
\thmref{c27-exchange} -- recognising the closed-neighbourhood superset that lets a greedy
stand in for the budgeted DP of \secref{c27-knapsack}.
\end{complexitybox}

\solstep{Pitfalls and Variants}
\begin{warningbox}
Ways to lose CF~1029E: (i)~attach at $v$ not $\parent(v)$ (Trap~1, over-counts on
brooms); (ii)~mark only $v$ when covering (Trap~2, siblings re-add edges);
(iii)~sort by index not depth (Trap~3, breaks post-order); (iv)~forget that
vertices of depth $\le 2$ are \emph{already} covered and add edges for them.
Variant: replace radius~$2$ with a general $R$ and attach at the $(R{-}1)$-th
ancestor -- the same exchange, one line changed. Contrast CF~1914F (Mixed Set): its
``pick a maximum antichain of non-ancestor pairs'' looks greedy but the correct
algorithm needs a path-invariant DP, so \thmref{c27-exchange}'s superset does
\emph{not} apply -- a live example of when \emph{not} to trust a greedy shadow.
\end{warningbox}

\subsection{Practice Problems}

\begin{stuckbox}
\begin{itemize}[nosep, leftmargin=1.4em]
  \item Did you prove (or at least test) the superset $\mathrm{cov}(\parent(v))
        \supseteq \mathrm{cov}(v)$ before committing to the greedy? If a counterexample
        exists, you need the DP of \secref{c27-knapsack}.
  \item Is your coverage mark the \emph{closed neighbourhood} of the attachment point,
        not just the deepest leaf?
  \item Does your sweep touch vertices in \emph{strictly decreasing depth} so each
        decision is final?
\end{itemize}
\end{stuckbox}

\begin{practiceset}
  \practice{Chloe and pleasant prizes (harvest two disjoint subtrees)}{Codeforces 743D}
    {https://codeforces.com/problemset/problem/743/D}{CF 1600 · L4}
    {Two players each cut one rope and take a whole subtree; the subtrees must be
    disjoint, and they want the maximum combined pleasantness. Root at~$1$, fold
    subtree sums, and at each node keep the \emph{best} subtree-sum available in
    each child branch (a greedy harvest, exactly the peak-combination of Ch.~13's
    \secref{c10-harvest}): the answer is the largest $m_1+m_2$ over a node's two
    child branches, and a chain makes it \texttt{Impossible}. Practise arguing the
    local ``keep the max seen so far'' choice is safe -- the same no-regret exchange
    as \thmref{c27-exchange}. Samples $25$, $2$, \texttt{Impossible} (hand-check
    the two-branch max at the root of the $8$-node example before submitting).}
\end{practiceset}

\paragraph{Back-cites (URLs live elsewhere in the book).}
\begin{itemize}[nosep,leftmargin=1.4em]
  \item \textbf{CF~763B ``Tree Queries''} (\chref{c10}'s mixed set owns this URL):
        a fixed-centre distance reasoning problem whose
        solution hinges on the same ``path to a fixed node'' structure -- contrast
        the \emph{decision} flavour there with the \emph{optimisation} greedy here.
  \item \textbf{Ch.~12 radius / \thmref{c09-dpjoin}}: the farthest-point-exchange
        behind Ch.~12's diameter is the sibling of \thmref{c27-exchange}; both lean on
        the unique path. Read \secref{c09-radius} for the centre-finding that fixes the
        vertex $1$ here for free.
\end{itemize}

\begin{recallbox}[10-second refresher -- 14.3 Greedy on trees]
A distance-from-a-centre cover collapses to a \emph{leaf-pruning greedy}
(\defref{c27-leaf-prune}): sweep by decreasing depth, and at each uncovered $v$ add
the edge to $\parent(v)$, marking $\parent(v)$'s closed neighbourhood. Optimal by
\thmref{c27-exchange} because $\mathrm{cov}(\parent(v)) \supseteq \mathrm{cov}(v)$ -- a
property of rooted trees (unique path), not general graphs. CF~1029E
(\probref{c27-greedy-solved}) is that greedy; on $T$ it adds $(1,9),(1,5)$, answer~$2$;
samples $2,0,1$.
\selftest{
  \item State the exchange in one sentence: why is attaching at $\parent(v)$ never
        worse than at $v$?
  \item Why does deepest-first make each decision final (no revisits)?
  \item What goes wrong if you mark only $v$ instead of the closed neighbourhood of
        the attachment point?
  \item Give one distance-cover instance where the greedy shadow is \emph{not} valid
        (costs vary) and say which section's DP you return to.
  \item Why is CF~1914F in the Mixed Set rather than solved here?
}
\end{recallbox}
